\documentclass[preview]{standalone}
\usepackage[english]{babel}
\usepackage{amsmath}
\usepackage{amssymb}
\usepackage{ragged2e}
\setlength\parindent{0pt}
\begin{document}
\begin{center}
\special{dvisvgm:raw <g id='unique000'>}\parbox{9cm}{\noindent{}\justifying{\special{dvisvgm:raw </g>}\special{dvisvgm:raw <g id='unique001'>}\begin{itemize} \item It is ``well known'' the groups acting an trees can be split into six mutually exclusives types [Ser03], [RS26].\end{itemize}\special{dvisvgm:raw </g>}\special{dvisvgm:raw <g id='unique002'>}\begin{itemize} \item Reid-Smith showed how these types can be detected from a local action diagram [RS26].\end{itemize}\special{dvisvgm:raw </g>}\special{dvisvgm:raw <g id='unique003'>}\begin{itemize} \item Chijoff-Tornier showed that the type of a group and the type of its (P)-closure are the same [CT26].\end{itemize}\special{dvisvgm:raw </g>}\special{dvisvgm:raw <g id='unique004'>}}}\special{dvisvgm:raw </g>}
\end{center}
\end{document}